\section{Discussion}
\label{sec:discussion}

\subsection{Elicitation}
The "first elicitation" session took noticeably more time than the others at 3 and a half hours. The following sessions all took no more than 2 and a half hours, showing that, as predicted in Sections~\ref{sec:value_functions} and~\ref{sec:qualitative_indicators_uncertainty}, the use of a first elicitation as a starting point for the following ones sped up the process considerably. An important aspect to discuss is the way the \gls{pilebwt} elicitation was carried out: in principle the comparison should have been reperated for each country beacuse of the intrinsic importance of the range of data for each indicator, which differ in this analysis, for example the "Electrcity amrket benefit" is a country-dependant indicator. However, to limit the required time the comparisons were always carried with a range that included the minimum and maximum data points across all countries. After the elicitation the elicited data point of a comparison was trasposed to the specific country range by means of linear mapping. This ensures that the country dependent value functions are still meaningul when applied to such data value, something that could not be possible if the transposition was to assume a non linear mapping, or if it was to just clip the values to the country range. To the author's knowledge, this approach should not introduce any significant bias in the results, however, further analysis could be carried out to confirm this assumption.

\subsection{Inconsistency Management}
Regarding the weighting process, expert E3 was extremely consistent in their judgments, helped by the fact that they were experienced in the use of \gls{mcda} methods. The other experts showed a varying degree of inconsistency, with expert E1 showing the highest inconsistency. As anticipated in \ref{sec:results}, expert E1 was very inconsistent in their \gls{pilebwt} judgments, however, the results obtained show that even including their input does not drastically change the final ranking, although, as expected, it does increase the uncertainty in the results.

Due to the nature of the \gls{pilebwt} method whcih has to be solved by means of optimization models, it was observed that inconsistencies in the judgments would prevent obtaining satisfactory results. Only cardinal consistency was considered in this analysis, that is, the order of the indicators cannot change between best comparisons and worst comparisons. For experts E2 and E4, some of these inconsistencies existed by a very low margin and were therefore corrected by slightly modifying the judgments. For example, the declared value by the expert in a comparison was 0.40, but to ensure consistency, it should have been greater than 0.45; the analysts adjusted the value to 0.45. For expert E1, some inconsistencies were much more severe and could not easily be corrected. However, the \gls{pilebwt} optimization model managed to run anyway, producing a wider weight space. To obtain the results of this thesis, both the results by excluding expert E1 and including them are shown, to demonstrate the effect of such inconsistent judgments on the final ranking.

\subsection{Observations on Expert Perspectives}
Thoughout the elicitations session it was clear that expert opinions were mostly aligned in their basis of judgment but ultimatly differed in the conclusions due to their different knowledge. This was expecially evident in the \glspl{qi} elicitations, where experts were more free to express their opinions. Context information in this regard proved to be the critial braking point between experts, for example, some experts would consider a certain licensing status as very good due to their knowledge of the technology, while those aware of political matters that affect some countries would rate the same status lower. This highlights the importance of selecting experts with diverse backgrounds to capture as many perspectives as possible. In the judgment of \glspl{smr} "insider" knwoledge was found to influence deeply the final outcome, were most would rate a technology at a certain level, the expert with direct experience with the same technolgy rated it lower due to their knowledge of recent issues that are not well known, even within the industry.

Regarding the \glspl{qi} themselves, the most critiqued one was the "design complexity" as many experts considered that it was possible to quantify it fairly easily by considering the number of safety systems, modularity level, redundancy level, etc. Second the "Construction Complexity" that the current study presented as country dependent was considered by most experts to be rather "country agnostic", as construction complexity would be mostly driven by the design itself rather than the country specific factors. Regarding "Construction Complexity", the author starting point was not preferring any design due to the belif that off-site construction and modularity do not bring a clear advantage in practice. All experts agreed on this viewpoint, mentioning that while modularity is, on paper, beneficial, it is yet to be proven and it is most likely to face several issues in their first of a kind projects.

All experts brough at some point in the conversation the significant influence that politics has on nuclear projects, behiond commitment, the ultimate choice of a reactor model is often driven by political factors and international relations rather than technical or economic merits. While this aspect is hard to quantify, it is important to acknowledge its role in the decision-making process for nuclear new builds.

Industry consultations revealed significant skepticism about SMR feasibility, primarily concerning vendor transparency. Experts questioned SMRs' market relevance, noting that most countries prioritize large-scale replacements for retiring capacity (e.g. replacement of coal-fired power plants in Poland). While the EPR's 1600 MWe capacity exceeds many nations' needs, the consensus identified 1000-1400 MWe as the optimal range for national deployment. SMRs were seen as potentially viable for private investors seeking smaller integrated solutions, but less suitable for utility-scale projects.

The proliferation of competing SMR designs was identified as problematic, creating decision paralysis for national entities while confusing public perception. Some observers suggested certain projects might represent "paper reactors" primarily aimed at securing funding rather than achieving commercial deployment.